\documentclass[a4paper,12pt]{article}
\usepackage[T1]{fontenc}
\usepackage[utf8]{inputenc}
\usepackage[polish]{babel}
\usepackage[margin=2cm]{geometry}
\usepackage{enumitem}
\usepackage{verbatim}
\usepackage{fvextra}
\usepackage{hyperref}
\RecustomVerbatimEnvironment{verbatim}{Verbatim}{breaklines=true,breakanywhere=true,fontsize=\small}

\title{Opis struktury projektu PlantYourHabit}
\date{10 listopada 2026}

\begin{document}
\maketitle

\section{Informacje ogólne}
PlantYourHabit jest aplikacją mobilną na system Android, napisaną w języku Kotlin. Interfejs użytkownika zbudowano z wykorzystaniem Jetpack Compose i biblioteki Material 3. Projekt składa się z jednego modułu aplikacji --- \texttt{app}. Główne funkcje aplikacji to ekran ogrodu z nawykami, lista questów oraz sklep z roślinami i motywami wizualnymi.

\section{Układ katalogów}
Poniższe drzewo przedstawia najważniejsze pliki projektu. Katalogi z zasobami zawierają dodatkowo pliki graficzne i konfiguracyjne, których nie wypisano pojedynczo.

\begin{verbatim}
PlantYourHabit/
|-- app/
|   |-- build.gradle.kts
|   `-- src/
|       |-- main/
|       |   |-- AndroidManifest.xml
|       |   |-- java/com/example/plantyourhabit/
|       |   |   |-- MainActivity.kt
|       |   |   |-- AppComponents.kt
|       |   |   |-- AppStateStore.kt
|       |   |   |-- AppTokens.kt
|       |   |   |-- GardenRoom.kt
|       |   |   |-- GardenScreen.kt
|       |   |   |-- GardenTheme.kt
|       |   |   |-- GardenTokens.kt
|       |   |   |-- GardenDialogs.kt
|       |   |   |-- Habit.kt
|       |   |   |-- PlantArtwork.kt
|       |   |   |-- PlantIntroVideo.kt
|       |   |   |-- PlantType.kt
|       |   |   |-- Quest.kt
|       |   |   |-- QuestsScreen.kt
|       |   |   |-- QuestsTokens.kt
|       |   |   |-- ShopOperations.kt
|       |   |   |-- ShopScreen.kt
|       |   |   |-- ShopTokens.kt
|       |   |   `-- ui/theme/
|       |   |       |-- Color.kt
|       |   |       |-- Theme.kt
|       |   |       `-- Type.kt
|       |   |-- res/
|       |   |   |-- drawable/       (ikony i ilustracje roślin)
|       |   |   |-- drawable-nodpi/ (grafiki okien i tła)
|       |   |   |-- mipmap-*/       (ikony aplikacji)
|       |   |   |-- raw/            (dźwięki i film wprowadzający)
|       |   |   |-- values/         (nazwy i motywy Androida)
|       |   |   `-- xml/            (reguły kopii zapasowej)
|       |   `-- keepRules/rules.keep
|       |-- test/java/com/example/plantyourhabit/
|       |   |-- ExampleUnitTest.kt
|       |   |-- HabitTest.kt
|       |   |-- QuestTest.kt
|       |   `-- ShopOperationsTest.kt
|       `-- androidTest/java/com/example/plantyourhabit/
|           |-- AppScreensComposeTest.kt
|           `-- ExampleInstrumentedTest.kt
|-- gradle/
|   |-- libs.versions.toml
|   |-- gradle-daemon-jvm.properties
|   `-- wrapper/
|-- build.gradle.kts
|-- settings.gradle.kts
|-- gradle.properties
|-- gradlew
`-- gradlew.bat
\end{verbatim}

\section{Moduł aplikacji}
\subsection{Punkt wejścia i przepływ danych}
\texttt{MainActivity.kt} uruchamia aplikację, konfiguruje motyw oraz wyświetla główny interfejs. W komponencie \texttt{PlantYourHabitApp} przechowywany jest stan ekranów, nawyków, questów, monet i elementów odblokowanych przez użytkownika. Wybór ekranu odbywa się przez dolny pasek nawigacyjny. Z tego miejsca przekazywane są też funkcje obsługujące działania użytkownika do poszczególnych ekranów.

\texttt{AppStateStore.kt} odpowiada za odczyt i zapis trwałego stanu aplikacji w \texttt{SharedPreferences}. Dane questów, monet, posiadanych roślin i motywów są kodowane do formatu JSON. Ustawienie jasnego lub ciemnego motywu jest przechowywane oddzielnie. Nawyki i przypisane im rośliny są obecnie przechowywane w stanie interfejsu; początkowa lista nawyków jest zdefiniowana w \texttt{Habit.kt}.

\paragraph{Odczyt i zapis stanu.}
\texttt{AppStateStore.load()} odczytuje JSON i odtwarza stan początkowy, jeśli zapis jeszcze nie istnieje. \texttt{save()} koduje stan do JSON i zapisuje go w preferencjach; niepowodzenie zapisu zgłaszane jest jako błąd. W \texttt{MainActivity.kt} zapis jest wywoływany ponownie, gdy zmieni się obiekt stanu przeznaczony do utrwalenia.

\begin{verbatim}
fun load(): PersistedAppState {
    val savedState = preferences.getString(STATE_KEY, null)
        ?: return PersistedAppState.initial()
    return decodeState(JSONObject(savedState))
}

fun save(state: PersistedAppState) {
    val saved = preferences.edit()
        .putString(STATE_KEY, encodeState(state).toString())
        .commit()
    check(saved) { "Failed to persist app state" }
}

LaunchedEffect(stateToPersist) {
    appStateStore.save(stateToPersist)
}
\end{verbatim}

\subsection{Interfejs użytkownika}
Ekrany i ich elementy są zaimplementowane jako funkcje \texttt{Composable}:
\begin{itemize}[leftmargin=1.2cm]
  \item \texttt{GardenScreen.kt} i \texttt{GardenRoom.kt} --- ekran ogrodu, obsługa nawyków oraz ich wizualizacja na półkach;
  \item \texttt{QuestsScreen.kt} --- lista questów, filtrowanie, dodawanie i edycja zadań;
  \item \texttt{ShopScreen.kt} --- sklep z roślinami i motywami ogrodu;
  \item \texttt{AppComponents.kt} i \texttt{GardenDialogs.kt} --- współdzielone elementy interfejsu i okna dialogowe;
  \item \texttt{PlantArtwork.kt} i \texttt{PlantIntroVideo.kt} --- grafiki roślin oraz ekran filmu wprowadzającego.
\end{itemize}

Pliki \texttt{AppTokens.kt}, \texttt{GardenTokens.kt}, \texttt{QuestsTokens.kt} i \texttt{ShopTokens.kt} grupują wartości używane do formatowania poszczególnych części interfejsu. Katalog \texttt{ui/theme} zawiera definicje kolorów, typografii i motywu Compose.

\subsection{Dane i reguły aplikacji}
\texttt{Habit.kt} zawiera model nawyku oraz reguły obliczania serii i etapów wzrostu rośliny. \texttt{Quest.kt} definiuje model questa i funkcje normalizowania tagów, filtrowania oraz przygotowania listy zadań. Typy roślin znajdują się w \texttt{PlantType.kt}, a dostępne motywy ogrodu w \texttt{GardenTheme.kt}. \texttt{ShopOperations.kt} zawiera operacje zakupu roślin i wyboru motywu wraz z kontrolą salda monet.

\section{Zasoby aplikacji}
Katalog \texttt{app/src/main/res} przechowuje zasoby niewykonywalne przez Kotlin. W \texttt{drawable} znajdują się ikony i ilustracje wykorzystywane w interfejsie, a \texttt{drawable-nodpi} zawiera grafiki tła i okien ogrodu. Katalogi \texttt{mipmap-*} zawierają ikony aplikacji w różnych gęstościach ekranu. W \texttt{raw} umieszczono dźwięki związane z nawykami oraz film wprowadzający. Pliki w \texttt{values} definiują m.in. nazwę i motyw systemowy aplikacji; pliki w \texttt{xml} określają zasady kopii zapasowej danych.

\subsection{Film wprowadzający}
Plik \texttt{app/src/main/res/raw/plant\_intro.mp4} jest odtwarzany przez \texttt{VideoView} osadzony w Compose za pomocą \texttt{AndroidView}. \texttt{MainActivity.kt} pokazuje ekran intro po uruchomieniu aplikacji i ukrywa go po zakończeniu filmu albo po wybraniu opcji pominięcia. Film jest wyciszony; w razie błędu odtwarzania aplikacja zapisuje błąd w logu i zamyka ekran intro. Po usunięciu widoku odtwarzanie jest zatrzymywane.

\begin{verbatim}
setOnCompletionListener { latestOnFinished.value() }
setOnErrorListener { _, what, extra ->
    Log.e("PlantYourHabit", "Intro video playback failed (what=$what, extra=$extra)")
    latestOnFinished.value()
    true
}
setOnPreparedListener { player -> player.setVolume(0f, 0f) }
setVideoURI(
    Uri.parse("android.resource://${context.packageName}/${R.raw.plant_intro}")
)
start()

DisposableEffect(videoView) {
    onDispose { videoView.stopPlayback() }
}
\end{verbatim}

\subsection{Dźwięki nawyków}
Dwa pliki dźwiękowe znajdują się w \texttt{app/src/main/res/raw}: \texttt{habit\_sound.wav} jest odtwarzany po oznaczeniu nawyku jako wykonanego, a \texttt{habit\_lose\_sound.wav} po cofnięciu wykonania. \texttt{MainActivity.kt} ładuje je przez \texttt{SoundPool}, czeka na zakończenie ładowania i odtwarza dźwięk tylko wtedy, gdy jest gotowy. Obiekt \texttt{SoundPool} jest zwalniany po usunięciu komponentu.

\begin{verbatim}
habitSoundId = habitSoundPool.load(context, R.raw.habit_sound, 1)
habitLoseSoundId = habitSoundPool.load(context, R.raw.habit_lose_sound, 1)

val done = today.toString() in habit.careDays
val soundLoaded = if (done) habitLoseSoundLoaded else habitSoundLoaded
if (soundLoaded) {
    val soundId = if (done) habitLoseSoundId else habitSoundId
    val streamId = habitSoundPool.play(soundId, 1f, 1f, 1, 0, 1f)
}
\end{verbatim}

\subsection{Grafiki roślin i ich formatowanie}
Ilustracje roślin znajdują się w \texttt{res/drawable} jako zasoby rysowalne Androida (m.in. pliki XML typu \texttt{vector}), a nie jako zdjęcia użytkownika. Przykładowe pliki to \texttt{ic\_plant\_fern.xml}, \texttt{ic\_plant\_monstera.xml}, \texttt{ic\_plant\_cactus.xml} oraz grafiki doniczek i kwiatów. \texttt{PlantArtwork.kt} wybiera ilustrację zależnie od gatunku oraz historii nawyku: na początku wyświetla nasiono, po przerwaniu historii pielęgnacji --- uschniętą roślinę, a w pozostałych przypadkach grafikę danego gatunku. Doniczka i opcjonalny kwiat są nakładane na roślinę jako osobne obrazy.

Obraz jest dopasowywany do całego obszaru rośliny przez \texttt{ContentScale.FillBounds}. Oznacza to wypełnienie dostępnej szerokości i wysokości, nawet jeśli proporcje grafiki różnią się od proporcji obszaru. \texttt{graphicsLayer} skaluje roślinę wraz z postępem nawyku, a filtr kolorów zmienia jej wygląd. Wzrost jest animowany, a dodatkowy kwiat pojawia się po osiągnięciu 75\% maksymalnej serii.

\begin{verbatim}
val plantModifier = Modifier
    .fillMaxSize()
    .graphicsLayer {
        scaleX = plantScale.scaleX
        scaleY = plantScale.scaleY
        transformOrigin = TransformOrigin(.5f, .78f)
    }

Image(
    painter = painterResource(resource),
    contentScale = ContentScale.FillBounds,
    colorFilter = if (state.animatedStage > 0f) state.colorFilter else null,
    modifier = state.modifier
)
\end{verbatim}

\section{Testy}
Testy jednostkowe znajdują się w \texttt{app/src/test}. Sprawdzają przede wszystkim obliczenia serii nawyków, logikę questów i operacje sklepu. Testy instrumentacyjne umieszczono w \texttt{app/src/androidTest}; obejmują testy ekranów Compose oraz test uruchomiony na urządzeniu lub emulatorze.

\section{Konfiguracja i budowanie}
Pliki \texttt{settings.gradle.kts} i główny \texttt{build.gradle.kts} konfigurują projekt Gradle i jego moduły. Konfiguracja modułu, identyfikator aplikacji, wersje SDK, opcje kompilacji i zależności znajdują się w \texttt{app/build.gradle.kts}. Katalog \texttt{gradle} zawiera katalog wersji bibliotek oraz pliki Gradle Wrapper, dzięki którym projekt można budować przy użyciu skryptu \texttt{gradlew} (Linux/macOS) lub \texttt{gradlew.bat} (Windows).

\end{document}
